\documentclass{article}

\usepackage{amsmath,amssymb}
\usepackage{xurl}
\usepackage[hidelinks]{hyperref}
\setlength{\emergencystretch}{2em}

% Shared commands for notes in docs/paper-gaps/.

\DeclareUnicodeCharacter{2080}{\ifmmode{}_{0}\else\textsubscript{0}\fi}
\DeclareUnicodeCharacter{2081}{\ifmmode{}_{1}\else\textsubscript{1}\fi}
\DeclareUnicodeCharacter{2082}{\ifmmode{}_{2}\else\textsubscript{2}\fi}
\DeclareUnicodeCharacter{2099}{\ifmmode{}_{n}\else\textsubscript{n}\fi}

\newcommand{\Lean}{\textsc{Lean}}
\ExplSyntaxOn
\NewDocumentCommand{\leanid}{m}{
  \ifmmode
    \textsf{\detokenize{#1}}
  \else
    \group_begin:
    \urlstyle{sf}
    \tl_set:Nn \l_tmpa_tl {#1}
    \tl_replace_all:Nnn \l_tmpa_tl {\_} {_}
    \exp_args:NV \path \l_tmpa_tl
    \group_end:
  \fi
}
\ExplSyntaxOff
\providecommand{\tr}{\operatorname{tr}}
\newcommand{\tnleanIssueBase}{https://github.com/LionSR/TNLean/issues/}
\newcommand{\tnleanPrBase}{https://github.com/LionSR/TNLean/pull/}
\newcommand{\tnleanDocBase}{https://sirui-lu.com/TNLean/docs/}
\newcommand{\ghissue}[1]{\href{\tnleanIssueBase#1}{GitHub issue \##1}}
\newcommand{\ghpr}[1]{\href{\tnleanPrBase#1}{PR \##1}}
\newcommand{\doclink}[1]{\href{\tnleanDocBase#1.html}{\path{#1}}}

% Dirac notation and the identity operator, as in blueprint/src/macros/common.tex.
% \providecommand keeps note-local definitions authoritative.
\usepackage{dsfont}
\providecommand{\ket}[1]{|#1\rangle}
\providecommand{\bra}[1]{\langle #1|}
\providecommand{\braket}[2]{\langle #1 | #2 \rangle}
\providecommand{\ketbra}[2]{|#1\rangle\!\langle #2|}
\providecommand{\Id}{\mathds{1}}
\providecommand{\one}{\mathds{1}}
\providecommand{\C}{\mathbb{C}}
\providecommand{\MN}[1]{M_{#1}(\C)}
\providecommand{\MD}{M_D(\C)}

% External referencing for paper sources.  The build helper writes sanitized
% auxiliary files under build/paper-aux/ and excludes duplicated source labels.
\usepackage{xr}

% Import a generated paper auxiliary file with a caller-supplied label prefix.
% The candidate paths support builds from docs/paper-gaps/, the repository root,
% and build/paper-aux/ itself.
\newcommand{\importpaperaux}[2]{%
  \IfFileExists{../../build/paper-aux/#2.aux}{%
    \externaldocument[#1]{../../build/paper-aux/#2}%
  }{%
    \IfFileExists{build/paper-aux/#2.aux}{%
      \externaldocument[#1]{build/paper-aux/#2}%
    }{%
      \IfFileExists{#2.aux}{%
        \externaldocument[#1]{#2}%
      }{}%
    }%
  }%
}

% General external reference macros
\newcommand{\paperref}[2]{\ref{#1-#2}}
\newcommand{\papereqref}[2]{(\ref{#1-#2})}

% CPSV16 source (1606.00608 / MPDO-22-12-17-2.tex) external document and shortcuts
\importpaperaux{cpsv16-}{1606.00608/MPDO-22-12-17-2-tnlean}
\newcommand{\cpsvref}[1]{\paperref{cpsv16}{#1}}
\newcommand{\cpsveqref}[1]{\papereqref{cpsv16}{#1}}

% Verdict marker: \gapnote{<kind>}{<status>}. Kind is the policy
% classification (clarification, local-correction, scope-restriction,
% unfaithful, false-source, open-gap); status is open, wip (elimination
% underway), resolved, or historical. Machine-read by the site index;
% prints a verdict line.
\newcommand{\gapnote}[2]{\par\noindent\textbf{Verdict:} #1 (#2).\par}

\providecommand{\leanid}[1]{\path{#1}}

\title{The nonempty-family condition in Wolf's Radon--Nikodym theorem}
\author{}
\date{2026-06-16}

\begin{document}
\maketitle

\section{Source statement}

Wolf's Radon--Nikodym theorem for quantum instruments states that, if
\(\{T_i\}\) is a family of completely positive maps satisfying
\(\sum_i T_i=T\) and
\[
  T(A)=V^\dagger(A\otimes \mathbf 1_r)V,
\]
then there are positive operators \(P_i\in\mathcal M_r\) such that
\[
  \sum_i P_i=\mathbf 1_r,
  \qquad
  T_i(A)=V^\dagger(A\otimes P_i)V.
\]
The local source is
\path{Notes/WolfNoteTexSource/ch02_representations.tex}:402--407, titled
``Radon--Nikodym for quantum instruments'' in
\cite{Wolf2012Quantum}.

\section{Local boundary}

The Lean theorem \leanid{IsCPMap.radon_nikodym_of_stinespring} makes the index
type nonempty.  This is not a strengthening of a mathematical hypothesis in the
substantive case; it is the finite-type boundary needed for the identity
\[
  \sum_i P_i=\mathbf 1_r
\]
to be meaningful as a resolution of the identity by instrument components.  If
the index type were empty, then the hypotheses \(\sum_i T_i=T\) force
\(T=0\).  One may still supply the zero Stinespring matrix
\(V=0\), but the conclusion would ask
\[
  0=\mathbf 1_r
\]
in \(\mathcal M_r\), which is false for a nonzero dilation dimension.

Thus the formal theorem states the source theorem in the nonempty instrument
case.  The blueprint entry records the same boundary explicitly by saying
``nonempty finite family.''

\section{Elimination plan}

No mathematical proof gap is expected for the nonempty case.  An empty-index
variant would need a different conclusion, for example a vacuous statement with
zero dilation space or no identity resolution.  Such a statement is not the
instrument theorem used in the development.

Verdict: local correction of a well-posedness boundary in the printed finite-family
statement.

\bibliographystyle{alpha}
\bibliography{references}

\end{document}
